% 190_T0_Korrekturen_De_ch.tex — Kompaktes Register (Neufassung 23. August 2026)
% Vollständige Begründungen: 190_T0_Korrekturen_Archiv_De

\section*{Vorbemerkung}

Dieses Dokument ist das kompakte, fortlaufende Korrekturregister des FFGFT-Korpus. Es ersetzt die bisherige ausf\"uhrliche Fassung, die unver\"andert als \textbf{Dok.~190-Archiv} (\texttt{190\_T0\_Korrekturen\_Archiv\_De/En}) vorliegt und f\"ur jeden Eintrag die vollst\"andige Begr\"undung, die fehlerhaften und korrekten Ausdr\"ucke im Wortlaut sowie alle Nachtr\"age enth\"alt. Die Tabelle unten f\"uhrt \emph{alle} Eintr\"age (K/C, P, R) mit den betroffenen Dokumenten und dem Gegenstand in Kurzform. Das Register bleibt append-only: Quelldokumente werden nicht revidiert (vgl.\ R50); neue Eintr\"age werden hier angeh\"angt, ihre Ausarbeitung erfolgt k\"unftig direkt in diesem Dokument.

\section{Registertabelle}

\noindent K/C~$=$~Korrektur, P~$=$~Pr\"azisierung, R~$=$~Revision/Registereintrag. Vollst\"andige Texte: Dok.~190-Archiv.

\begingroup\small
\begin{longtable}{@{}p{0.9cm}p{3.1cm}p{10.2cm}@{}}
\toprule
\textbf{Nr.} & \textbf{Betrifft} & \textbf{Gegenstand (Kurzform; Ausf\"uhrung: Archiv)} \\
\midrule\endhead
\bottomrule\endfoot
K1 & 049 (HTML) & $\xi = \lambda_h^2 v^2/(16\pi^3 m_h^2)$ \emph{[Nachtrag 28.~Juli 2026, vgl.\ R66: Grund der Korrektur --- $16\pi^3 = (4\pi)\cdot(4\pi^2) = S^2\cdot 2S^3$ ist die Sphärenzählung der kompakten T$^4$-Randgeometrie (Dok.~310); die frühere Fassung zählte eine 2-Sphäre zu viel]} \\
K2 & 116 (Tab.) & $r_\tau = 25/9$ \\
K3 & 018 (Explorer) & $a_e = 2\pi^2/(f \cdot k)$ \\
K4 & 267 (DE+EN), 268 (DE) & \textbf{Faktor~3 $=$ drei beobachtbare Raumdimensionen:} der Beobachter misst $k_\mathrm{obs}^2=(n_1^2+n_2^2+n_3^2)/L^2$ (die vierte Torusrichtung ist Zeitentfaltung). \\
K5 & 275 (Skript v3) & \textbf{Alle drei behoben in v4} (11.~Juni 2026): $T^4$ als kNN in echten 4D-Koordinaten mit dimensions-gematchten 4D-Nulls (Ergebnis: $1{,}329\pm0{,}025$ \emph{im} Band $[0{,}690,\,2{,}406]$); HLV-Ensemble mit $\sigma=0{,}05$ ($0{,}710\pm0{,}008$, im Band); Schalen-Detektor wirf\,\ldots \\
P1 & 116 / 158 & Beide korrekt; $0{,}001\,\%$ für extern \\
P2 & mehrere & $\hbar$ folgt aus $\xi$ via $L_0 = \xi\cdot\ell_P$ \\
P3 & 173--176 & Zwei Parameter: $\xi_0$ (geom.) und $\xi_{\text{Higgs}}$ (alg.) \\
P4 & 180--182 & $L_0$ = geom.\ Fundamentallänge \\
P5 & mehrere & Massenstruktur aus $\xi$; $\sim 1\,\%$ Übereinstimmung \\
P6 & 116 / 158 & Koide nur mit numerischen Werten (Nichtabschluss, Dok.~193) \\
P7 & 005, 008, 019, 086, 189, 202 & Skalenstruktur aus der Rekursion $\mathcal{R}$ \\
P8 & 005, 008, 012, 014, 041, 060, 133, 141, 159, 181 & Geometrische Skalen\-anpassung / fraktale Korrektur \\
P9 & 005, 006, 042, 046 & Spalten \enquote{Skalierungs\-klasse} + \enquote{SM-Korrespondenz}; \newline Quarks via $m = r\,\xi^p\,v$ \\
P10 & 009 & $K = R_{pe}/(4/3)^7$: ein empirischer Anker; \newline Faktor $1{,}314$ implizit definiert, nicht hergeleitet \\
P11 & 025, 003, 133 & Korrekt: $(1-275/4\cdot\xi)$, $\Delta=0{,}0002\,\%$; \newline Dok.~003 und Dok.~133 verwenden verschiedene Modelle \\
P12 & 022, 230 & Zwei verschiedene Observablen: kumulativ (Dok.~022) vs.\ elementar $\xi/(2\pi)$ (Dok.~230); $10^{-5}$ bei keinem endlichen $N$ aus Dok.~022 erreichbar \\
P13 & 013, 025, 061 & $\xi_\text{FFGFT}/\xi_\text{UIFT} = m_ec^2/(\frac{16}{9}\ln 2\cdot\text{eV}) \approx 414\,684$; \newline Tabelle in Dok.~013 (archiviert) rekonstruiert \\
P14a & 041 & Fraktale Wegverlängerung (geometrische Wegstreckung); Energieverlust nur Artefakt bei weggelassener Wegverlängerung; Weglängen-Abhängigkeit bleibt (Dok.~026/149/182); dieselbe Wegverlängerung trägt die CMB-Behandlung in Dok.~268 (vormals separat als P19 geführt) \\
P15 & 061 / 041 & Als $\Lambda$CDM-Vergleichsgröße markiert; statisches Universum kennt kein kosmologisches $z$; CMB durch stationäre Thermalisierung \\
P16 & 026 / 064 / 181 & $H_0$ ist emergent (dekompaktifiziert), nicht fundamental; in Dok.~263 umgesetzt, Dok.~026/064/181 als veraltet vorgemerkt \\
P17 & 028 & DE/DM-Relation ist zirkulär (ergibt 2,5 für jedes $\xi$); MOND teil-hergeleitet ($\xi^{1/4}$ echt, $K_M$ frei); CMB-Relationen unberührt; vorgemerkt \\
P18 & 267 & Beide Bilder sind über die beobachteten Verhältnisse \emph{nicht} unterscheidbar; die Entscheidung ist theoretisch, nicht empirisch (konzeptuelle Einordnung) \\
P20 & 267 / 268 & Die \emph{Form} liegt fest: dimensionsloses $\xi$-Verhältnis ($R_H/\ell_P$). \\
P29 & 268 & \emph{Nur teilweise beantwortet.} Der Faktor~3 aus der 3D-Projektion ($3\cdot\{1,6,14,26\}$) ist plausibel, aber die Folge $\{1,6,14,26\}$ ist \emph{aus den Daten abgelesen} (Dok.~268, Schritt~4), nicht vorwärts erzeugt. \\
P30 & 268 & Die \emph{naive} sphärische Bessel-Summe $C_\ell=\sum_n N_\mathrm{BE}|j_\ell(k_\mathrm{obs}D_A)|^2$ reproduziert $\{3,18,42,78\}$ \emph{nicht} -- sie wäscht zu einem Quasi-Kontinuum aus, dominiert von niedrigen Moden. \\
P31 & 268 & \emph{Führungsverhalten forward hergeleitet.} Forward gibt T$^4$ das harmonische Rückgrat $1:2:3:4:5$ (symmetrische Resonanz $k^2=3n^2$). \\
P32 & korpusweit & Folgt aus P20: $H_0$ ist eine \emph{importierte} externe Kalibrierung, nicht hergeleitet. \\
P33 & 009 / 042 / 181 & Die \emph{Form} von $\xi$ ist vorwärts: das harmonische $4/3$ (reine Quarte, Dok.~042/159) gibt die Leitziffer, $3$ (Raumdimensionen) und $2^2{=}4$ (Raumzeit) liest man aus der Geometrie. \\
P34 & 181 / 013 & \emph{3+1 ist empirische Eingabe} -- keine Theorie leitet die Dimensionszahl her; der frühere \enquote{folgt zwingend}-Anspruch (Dok.~181) ist überstellt. \\
P35 & korpusweit (bes.\ 182, 250er, Skalentabellen) & \textbf{$\xi^N$-Trivialität:} \enquote{$X = \xi^N$} ist für sich \emph{aussagelos}, da sich jede Magnitude als $\xi$-Potenz schreiben lässt -- ein Exponent allein ist keine Vorhersage (vgl.\ P10, P20, P33: Form vorwärts, Zahl Eingabe). \\
P36 & 182 / kosmolog.\ Sektor & \textbf{$R_H$-Deutung präzisiert:} $R_H = c/H_0$ ist die \emph{Hubble-Länge} -- die lokale Expansionsrate als Distanz ausgedrückt -- \emph{nicht} die Universumsgröße und \emph{nicht} der beobachtbare Horizont (dieser $\approx 3\,R_H$ in $\Lambda$CDM). \\
P37 & 009, 188, 271, 272, 275 & \textbf{Drei Ebenen klar getrennt:} \newline \textbf{(1) Mikroskopisch -- Input:} $\xi = 4/30000$ (fraktale Dimension). \\
P38 & 275 & \textbf{Statusklärung:} (a)~$1/\varphi$ ist \emph{Durchgangspunkt}, kein Fixpunkt -- der einzige Fixpunkt von $r\mapsto r(1-\xi)$ ist $0$. \\
P39 & 182, 267, 277, korpusweit ($H_0$/$R_H$) & \textbf{$H_0$/$R_H$-Status präzisiert (gemeinsamer Bezugspunkt, keine FFGFT-Lücke):} Die Relation Plancklänge$\leftrightarrow R_H$ (Exponent $41/4$, $R_H/L_0 = 6{,}49\times10^{64}$) wird von \emph{keinem} Rahmen aus ersten Prinzipien hergeleitet -- warum das Universum $\sim$\,$10^{60}$\,\ldots \\
K6 & 004, 025, 030, 039, 041, 053, 061, 063, 068, 081 & \textbf{Rotverschiebung ist achromatisch.} Die fraktale Wegverlängerung ist rein geometrisch und streckt \emph{alle} Wellenlängen mit demselben Faktor: $1+z=e^{\xi x}$, unabhängig von $\lambda$. \\
P40 & 005, 006, 011, 012, 133, 275 & \textbf{Nur Verhältnisse sind exakt.} Absolutwerte tragen die fraktale/SI-Korrektur, absorbiert in die Brückenkonstanten $v$ (Massen $m=r\xi^p v$), $E_0=7{,}398$~MeV ($\alpha$) und $3{,}521\times10^{-2}$/$2{,}843\times10^{-5}$ ($G$; \texttt{calc\_De.py}); in Verhältnissen kürzen\,\ldots \\
P41 & 267, 280, 221 (kosmolog.\ Sektor) & \textbf{Geprüft und nicht gebucht (Entartungs-Risiko an der geteilten Skala).} $\Lambda$CDM-Signal $\Delta v=c\,\dot z\,\Delta t/(1+z)$ über 25 Jahre: $+6$\,cm/s ($z{=}0{,}5$), Vorzeichenwechsel bei $z_0=1{,}92$, $-14$\,cm/s ($z{=}4$), $-20$\,cm/s ($z{=}5$); ELT/ANDES-Größenordnu\,\ldots \\
P42 & 282, 292 (Leptonsektor) & \textbf{Ein deklarierter Referenzpunkt: das Verhältnis $m_\mu/m_e$.} Der Rahmen bleibt grundsätzlich bei der rein theoretischen Ableitung aus $\xi$; der Leptonsektor erhält -- analog zu P39 ($H_0$/$R_H$ als gemeinsamer empirischer Bezugspunkt) -- genau \emph{einen} deklarierten R\,\ldots \\
P43 & 006, 011, 292 (\texttt{calc\_De.py}) & \textbf{Kein Umbau; als Notiz gebucht.} Nachgeprüft: \texttt{calc\_De.py} verwendet $K_{\text{frak}}=1-100\xi$ \emph{nicht explizit}. \\
R41 & 284 (vgl.\ 270, 249) & \textbf{Der Verlust ist nicht die SI-Umwandlung.} Nach Dok.~270 ist die Typ-I-Dualitäts-Dekompaktifizierung $S^1_m\!\to\!\mathbb{R}_t$ ($\tilde{T}\cdot m=1$) ein \emph{mode-erhaltendes Embedding} (verlustfrei); die Kompaktheit überlebt als diskretes Spektrum und als physikalische\,\ldots \\
R42 & 270 (vgl.\ 248) & \textbf{Keine offene geometrische Frage -- es ist der Messvorgang.} Im kompakten Bild sind die vier Kreise symmetrisch (keine Raum/Zeit-Unterscheidung); die Zeit wird nicht von der Geometrie gewählt, sondern ist \emph{relational} -- die Symmetrie-Brechung \emph{ist} der Messvorga\,\ldots \\
R43 & 285, 283 & \textbf{Auf Enthaltensein korrigiert.} $C_3<A_5$ ($3$- und $5$-Zähligkeit permutieren denselben Ikosaeder, \texttt{recompute\_reconcile.py}) und $T^7=T^4\times T^3$ ($7=4+3$), also $\mathrm{HLV}\supset\mathrm{FFGFT}$ auf Gruppen-/Dimensionsebene -- die Modelle sind vereinbar, nic\,\ldots \\
R44 & 272, 190 (P35), 285 & \textbf{Pr\"azisiert (kein Widerspruch).} $\varphi$ erh\"alt Bedeutung \emph{nicht} in FFGFT, sondern als hergeleiteter Eigenwert von HLVs $5$-fach-/Perp-Sektor (golden inflation $\varphi^2=\varphi+1$) -- Marker der Grenze HLV$\setminus$FFGFT. \\
R45 & 285, 282--284 & \textbf{Pr\"azisiert.} Ein Quasikristall hat zwei Spektren: das \emph{Struktur-/Beugungsspektrum} ist rein punktf\"ormig (scharfe Bragg-Peaks, $\approx21\times$ \"uber Zufalls-Nullmodell, \texttt{quasicrystal\_information.py}) -- das ist die Information und genau HLVs Substrat-Un\,\ldots \\
R46 & 296 (vgl.\ 025, 026, 063, 156) & \textbf{Offener Prüfpunkt (noch nicht bearbeitet).} Diese Zwischenskalen-Formulierung steht potenziell \emph{quer} zur bereits etablierten Korpus-Linie, die Dunkle Materie/Energie nicht als offene Signatur, sondern als \emph{abgeleitete Konsequenz} behandelt: Dok.~025 deutet Dunk\,\ldots \\
R47 & 285, 297, 298 (vgl.\ 286, 284, 295) & \textbf{Im Zeitsektor: Enthaltensein, nicht Vergrößerung.} In einheitlicher Schreibweise -- beide $D6+$Zeit, Marcel als $D6+$Spiralzeit, FFGFT im Hilbertraum als $D6+$Rekursionszeit -- ist das Spiralzeit-Fenster nicht additiv, sondern \emph{intern}: ein beschränkter Ausschnitt (F\,\ldots \\
R48 & 304 (vgl.\ 295, 301, 303) & \textbf{Notiz zurückgenommen, Dokument präzisiert (nicht repariert).} Die Überzeichnung lag allein in der \emph{Mail} und wurde ohne Vorbehalt zurückgenommen; das Dokument schrieb durchgehend aus der FFGFT-Lesart und buchte die Brücke bereits als restricted. \\
R49 & 304, 305 (vgl.\ 295, 303, 272) & \textbf{Vokabel zitiert, als Analyse markiert, Abstract entkürzelt --- Notation nicht umbenannt.} Eine konsolidierte Erstnennungs-Fußnote attribuiert U1/U2/U3 \emph{und} $M_\tau/M_B$ an Krüger (2026) mit Primärquelle (\texttt{doi:10.5281/zenodo.21302278}) und dokumentierter Linea\,\ldots \\
R50 & 025, 063, 061, 064, 078 (ersetzt durch 306; vgl.\ 003, 010, 295, 298, 302) & \textbf{Begründung ersetzt, Konklusion unverändert.} Die Konklusion (kein singulärer Anfang; endlicher Kern mit minimaler, von $\xi$ gesetzter Skala $L_0$) bleibt. \\
R51 & 049 (vgl.\ 095, 306, 267, 302; Notation) & \textbf{Notations-Altlast, kein Substanzwiderspruch --- und dennoch nicht mehr maßgeblich.} \emph{(a) Die Auflösung:} In der ursprünglichen T0-Sicht bezeichnete $T_0$ \emph{nie} den Zeitpunkt null, sondern den \emph{kleinsten möglichen Zeitraum}. \\
R52 & 101, 165 (Zahlwerte; vgl.\ 180, 182, 013, 250, 290) & \textbf{Zahlwerte korrigiert; Hauptrelation unberührt.} \emph{(i) Falscher $L_0$-Wert:} korrekt ist $L_0 = \xi\,\ell_P = \tfrac{4}{30000}\cdot 1{,}616\times10^{-35}\,\text{m} = \boxed{2{,}155\times10^{-39}\,\text{m}}$ --- so in Dok.~180 (der eigens dafür bestimmten Herleitung), D\,\ldots \\
R53 & 267 \S518 (Tabellenzeile \emph{Quelle der Skala}; vgl.\ 182, 165, P20/P35/P36) & \textbf{Korrigierte Lesart --- die Symmetrie wird dadurch schärfer, nicht schwächer.} \emph{(a)} Die Tabellenzeile \emph{Quelle der Skala} lautet auf FFGFT-Seite korrekt nicht \enquote{aus $\xi$}, sondern \textbf{\enquote{gesetzt / an $H_0$ kalibriert}}. \\
R54 & 292, 306 (vgl.\ 307) & \textbf{(a) $N_0$ ist keine Grundeinheit.} $N_0\approx38{,}6$ ist ein \emph{leiter-interner}, generationszählender Faktor ($N_g=g\,N_0$, Dok.~292); sein Wert ist \emph{offen} (P35: $100/e$, $100/\varphi^2$, $39$ treffen ihn nicht sicher). \\
R55 & 292 \S{}Leiterübergreifende Kopplung (vgl.\ R54, P35, P40/P41) & \textbf{Das Gesetz bleibt bestehen; seine Evidenzbreite ist eine Stelle, nicht drei.} \emph{(i) Die dritte Stelle ist arithmetisch erzwungen.} Aus $m_\tau/m_e=(m_\tau/m_\mu)(m_\mu/m_e)$ folgt für eine durchgängig multiplikative Korrektur $K^{p_3}=K^{p_1}K^{p_2}$, also $p_3=p_1+p_2$\,\ldots \\
R56 & 188 \S{}Besondere Rolle von $n=3$ (vgl.\ 189 \S{}Ableitungshierarchie, 248, 230/231/232, R55) & \textbf{Die drei Elemente werden hiermit einheitlich als Grundannahmen gebucht.} \emph{(i) Der deklarierte Rahmen.} FFGFT setzt drei Strukturelemente: die Grundrelation $\tilde{T}\cdot m = 1$ (Dok.~189, Stufe~1), die Topologie $T^4$ (Dok.~248) und die Ordnung $Z_3$ der Identifika\,\ldots \\
R57 & A015 (vgl.\ 020, 030, A-Serie) & \textbf{Forderung aufgelöst --- $\xi$ ist Fundamentalparameter.} Ein Beweis ist weder möglich noch nötig; die Forderung war falsch gestellt. \\
R58 & A142, A140 (vgl.\ 127, 163) & \textbf{Gravitationsquantisierung nicht gefordert.} $G=\xi^2/(4m_\text{char})$ ist intrinsisch aus der Torus-Geometrie; eine separate Quantengravitation entsteht nicht, weil die Geometrie $T^4/Z_3$ pre-geordnet und nicht dynamisch im selben Sinne wie Felder der Quantenfeldtheorie\,\ldots \\
R59 & A130 (vgl.\ 011, 043, 101; P40) & \textbf{Strukturell erklärt [S].} Zwei Ursachen: (i)~SI-Werte sind schemaabhängig (1-Loop, $\overline{\mathrm{MS}}$), d.\,h.\ die EFT-Präzisionswerte sind keine theorie-unabhängigen Rohdaten; (ii)~fraktale Korrekturen in $v$ sind nicht sauber von der Schleifenstruktur des SM tren\,\ldots \\
R60 & A160, A165 (vgl.\ 020, 022, 023, 230; P35) & \textbf{Beide Faktoren jetzt geometrisch hergeleitet [B].} (i)~$2\pi$ ist der Kreisumfang des Einheitskreises --- die natürliche Umrechnung zwischen der Torus-Windungsphase und dem Messwinkel; $\xi/(2\pi)$ ist damit die Phasenkorrektur pro Messung. \\
R61 & 174, 175 (vgl.\ A142, 147, 162; R57, P35) & \textbf{Datengetriebene Wertauswahl --- kein Zweitparameter; zurückgezogen.} \emph{(i) Der Hergang:} Die Auswahl von $\xi_\text{Higgs}$ erfolgte, \emph{weil} er das gemessene Fenster trifft; die Zwei-Räume-Doktrin ist Rationalisierung nach der Auswahl, keine unabhängige Begründun\,\ldots \\
R62 & 005 \S{}Baryon-Beispielrechnung (Proton) (vgl.\ R50, P35) & \textbf{Der Anker der Formelzeile ist $\pi^2/2$, nicht $m_\mu$.} \emph{(i) Diagnose.} Der Fehlfaktor identifiziert den Fehler als ein einziges Symbol: $(\pi^2/2)/m_\mu = 46{,}705$ trifft den Fehlfaktor auf $0{,}07\,\%$. \\
R63 & 253 \S{}$\xi$-Resonanzheuristik (Code: \texttt{rsa/\allowbreak factorization\_\allowbreak benchmark\_\allowbreak library.py}; vgl.\ 024, 075, 076, 176; R57, P35) & \textbf{$\xi$ ist im \enquote{optimalen} Bereich wirkungslos; die Optimierung hat den Filter abgeschaltet.} \emph{(i) Inertheit.} Der Schwellwert des Codes ist $1/1000$. \\
P44 & 253, 190 (R63); vgl.\ A135, P35 & \textbf{Vorregistriert geprüft und für diesen Aufbau widerlegt; die Fensterbreite ist eine Eigenschaft des Messaufbaus.} \emph{(i) Aufbau.} Shor-Phasenschätzung, $12$ Zählqubits (Phasenauflösung $2^{-12} = 2{,}44\times10^{-4}$), $8$ Schüsse, $10\,\%$ vollständig zufällige Messerg\,\ldots \\
R64 & \texttt{2/python/t0\_no\_fallback\_shore.py}; & \textbf{Das Verfahren ist richtig, die Trennebene zu grob und das Etikett falsch.} \emph{(i) Was zutrifft.} Die Zähler \texttt{t0\_success\_count} und \texttt{fallback\_success\_count} werden an getrennten Stellen erhöht (Z.~376 bzw.\ 280/300), \texttt{get\_statistics()} gibt bei\,\ldots \\
R65 & 024, 075, 076, 147, 176, 253 sowie der & \textbf{Es existiert kein von Shor unterscheidbares Verfahren; die Methodenbehauptung wird zurückgezogen.} \emph{(i) Was zurückgezogen wird.} (a)~Die Behauptung eines eigenständigen T0- bzw.\ $\xi$-Faktorisierungsverfahrens. \\
R66 & 097 (DE/EN sowie die Standalone-Fassungen unter & \textbf{Die $16\pi^3$-Zählweise ist richtig --- aber nicht aus dem in Dok.~097 angegebenen Grund.} \\
K7 & die vierzehn HTML-Seiten unter \texttt{rsa/} & \textbf{Umgebaut, nicht mit Hinweisen versehen: die Seiten selbst sagen jetzt, was sie tun.} Kein Korrekturkasten wurde eingefügt; korrigiert wurden Rechenweg, Beschriftung und Methodentext. \\
R67 & A040 (fraktale Korrektur), A270 (Bulk-Exponent 36), & \emph{(i)}~Bei $D_f^\text{eff} = 2{,}973$ (drei Nachkommastellen) lässt das Rundungsintervall $n \in [98{,}3,\ 102{,}0]$ offen; selbst $D = 2{,}973$ exakt ergibt $n = 100{,}12$. \\
R68 & 279 ($H_0$-Kette), 308 ($\Lambda$-Lesart, & \emph{(i)~Identifikation:} Die $\Lambda$-Funktion wird von der abgeleiteten Größe $\Lambda^{*} := 1/R_H^2 = (\pi/2)^2\xi^{20}/ \bar\lambda_e^2 = 5{,}218\times10^{-53}$~m$^{-2}$ getragen (\enquote{Abschlussskala}; Namensdisziplin analog $a_0$). \\
R69 & 308 \S{}$a_0$-Herleitung; vgl.\ P20, R68; & \emph{(i)}~Aus $\partial\text{T}^4 = \emptyset$ (Dok.~306) folgt bei isotropem Torus ein kompakter Zeitzyklus $\tau_c = L^{*}/c = 2\pi/H_0 = 91{,}9$~Gyr [K$|$P20]; das Energiequant ist exakt $\hbar H_0$ und identisch mit dem Impulsquant der Raumzyklen. \\
R70 & 061 (Temperatur/CMB), A085 & $T_0$-Skalenlücke \\
R71 & 309 (Skalenanker, Vererbungsproblem); & Verschärfung: Seitenwahl in der
Hubble-Spannung \\
R72 & A010 (Glossar), A080 (Einheiten), A265 & $E_0$: nackt
und korrigiert \\
R73 & 026 (Geometrische Kosmologie), 279 & $H_0$-Schreibweisen, Exponent
$41/4$ \\
R74 & 025 (Kosmologie), 063 (Cosmic), jeweils & Lithiumproblem: Mechanismus
statt Lösung \\
R75 & 024, 034, 075, 076, 147, 176: Namenskollision $\xi$
und Neufassung der Shor-Dokumente & --- \\
R76 & 318 & Geltungsbereich der Masseableitungen \\
R77 & 041, 160, 318 & Einordnung früher Kopplungskonstanten-Formeln \\
R78 & 160, 005 & $\alpha_s$ und laufende Kopplung \\
R79 & 321, 319, 049 & Algebraische Herleitung der $SU(3)_c$-Eichgruppe \\
R81 & 323, 321, 160, 320 & Weinberg-Winkel bei $M_Z$ aus $\xi$ \\
R82 & 322, 320 & Hilbertraum-Einbettung und Spektraltheorie \\
R83 & 319, 321, 323, 322, 320 & Statusaktualisierung geschlossener Brücken \\
R84 & 324, 321, 009 & Casimir-Herleitung von $\xi$; $\xi$ kein G(6)-Spektralwert \\
R85 & 313, 325, 257, 321, 322 & Hawking-Mechanismus mikroskopisch geschlossen \\
R86 & Dok.~180, 322, 314, 321, 325 & Selbstadjungiertheit $\hat{F}$ vollst. geschlossen; $\xi = \lambda_{\min}(\hat{F}_{D_4})$ \textbf{[B]} \\
R87 & Dok.~320, 322 & $\Delta m_{32}^2$: Zahlenwerte korrigiert, Abw. $+16{,}0\,\%$; Kandidatenanalyse erg\"anzt \textbf{[K]} \\
R88 & Dok.~326, 329 & $n_{\mathrm{thresh}}$ skalenabhängig: $n(L)=L/(\sqrt{2}\ell_P)$; Matzkes Skala $=2\pi\ell_P/\ln 2$ \textbf{[B]}, universelle Zahl \textbf{[X]} \\
R89 & Dok.~180, 019, 201, 032 & $m_T = M_\mathrm{Basis}/\xi$: eine Struktur, sektorabh\"angige Basismasse; Reichweite $=\xi\cdot\lambda_C(M)$ \textbf{[B]} \\
R90 & Dok.~180, 019, 201 & Exponent in $L_0 = \xi^1\ell_P$ aus $m_T\sim1/\xi$ hergeleitet, nicht analogisiert \textbf{[B]} \\
R91 & Dok.~019, 201 & $\lambda$ in $m_T=\lambda/\xi$ ist eine Basismasse, nicht der Higgs-Quartic --- Beschriftung irref\"uhrend \textbf{[K]} \\
R92 & Dok.~250, 190 & Dok.~250: $\xi/\ell$ ist Kopplung $g_T$, nicht $m_T$; Schwarzschild-Probe f\"ur $L_0$ ist Identit\"at \textbf{[B]} \\
\end{longtable}
\endgroup

\section{Aktuell offene Br\"ucken (Stand R85)}

\begin{itemize}\itemsep2pt
  \item $\Delta m_{32}^2$ atmosph\"arisch ($+16{,}0\,\%$); Mischungsterm $F_5$--$F_7$ --- Dok.~320, R87 [S]
  \item 2-loop-Korrekturen zu $\alpha_s(M_Z)$ [S]
  \item $m_{\mathrm{Pl}}$ und $\alpha_{\mathrm{em}}(M_Z)$ aus $\xi$ (externe Eingaben in R81) [S]
  \item Quark-/Hadronsektor --- Dok.~318, R76
  \item Grauk\"orper-Faktoren; R\"uckreaktion auf Orbifold-Fixpunkte --- R85 [S]
  \item Kosmischer Exponent $41/4$ vorw\"arts herleiten (P20); daran h\"angen die $10^{123}$-Frage, $R_H\to\ell_1$ und $\kappa = \Lambda_{\mathrm{obs}}R_H^2$
  \item CMB-Peak-Selektion: Folge $\{1,6,14,26\}$ vorw\"arts erzeugen; $|n|^2{=}30$-Fall (P29/P31)
  \item $C_\ell$-Projektion: physikalische Quell-/Fensterfunktion (P30)
  \item $\Omega_{\mathrm{DM}}$ quantitativ modellneutral (P17)
  \item Spektraldimension $d_s = 1{,}86$ vs.\ FFGFT $\approx 2$
  \item $H_0$-Sprachbereinigung korpusweit (P34); Dok.~026 grundlegend \"uberarbeiten (P16)
  \item Basismasse $M_\mathrm{Basis} = m_P$ im fundamentalen Sektor begründen --- Dok.~180, R90 [S]
  \item $\xi$ als Linienbreite (P44, vorregistriert); $\iota$-Einbettung HLV$\,\subset\,$FFGFT im Zeitsektor (Dok.~297)
\end{itemize}
\textbf{R86 -- Selbstadjungiertheit von $\hat{F}$ (August 2026):}
Dok.~327 schließt die in R82 als [S] deklarierte Lücke vollständig.
Die Endlichkeit der Skalenleiter ist durch $L_0 = \xi\cdot\ell_P$
(Dok.~180 \textbf{[K]}) garantiert -- keine Unendlichkeiten in FFGFT,
damit $\hat{F}$ beschränkt und $\mathcal{D}(\hat{F}) = \mathcal{H}$.
Die Symmetrie $\hat{F} = \hat{F}^\dagger$ folgt aus der $\mathbb{Z}_3$-Paarung
$(k,-k)$ (dieselbe wie in Dok.~325 \textbf{[B]}).
Defektindizes $(0,0)$: genau eine selbstadjungierte Realisierung \textbf{[B]}.
Fraktales Maß (endliche 100er-Rekursion), $\mathbb{Z}_3$-Restriktion
und alle drei $\chi$-Twist-Klassen erhalten die Selbstadjungiertheit \textbf{[B]}.
$\xi = \lambda_{\min}(\hat{F}_{D_4})$ wohldefiniert und trunkierungsstabil
(Min-Max) \textbf{[B]}.
Es verbleiben keine Restfälle.
Prüfskript: \texttt{327\_selbstadjungiertheit\_F\_verify.py} (21 Assertionen).\\
\emph{Deklariert 23.~August 2026}

\medskip

\textbf{R87 -- $\Delta m_{32}^2$: korrigierte Zahlenwerte und Kandidatenanalyse (August 2026):}
In Dok.~320 war die Auswertung von $\xi^{9/5}$ fehlerhaft ($8{,}717\times10^{-8}$ statt
$1{,}059\times10^{-7}$); die Massenformel $m_{\nu_i} = m_e\,\xi^{p_i}$ und die Exponenten
$p_i \in \{9/4,\,2,\,9/5\}$ waren stets korrekt. Dok.~320 und Dok.~322 wurden auf die
richtigen Zahlenwerte gebracht: $m_{\nu_3} = 54{,}11$~meV \textbf{[K]},
$\Delta m_{32}^2 = 2{,}846\times10^{-3}$~eV² \textbf{[K]} (Abweichung $+16{,}0\,\%$ gegen
NuFIT~5.3), $\sum m_\nu = 64{,}17$~meV \textbf{[K]} (weiterhin im Planck-Limit).
Unberührt: $m_{\nu_1}$, $m_{\nu_2}$, $\Delta m_{21}^2$ ($+8{,}3\,\%$) und die
Fixpunkt-Zuordnung. Dok.~320 wurde zugleich um eine Kandidatenanalyse für die
verbleibende Abweichung erweitert: $K_{\mathrm{frak}}$ am Fixpunkt $F_7$ ($+12{,}9\,\%$,
unzureichend), Mischungsterm $F_5$--$F_7$ (richtiges Vorzeichen, vollständige
Massenmatrix steht aus -- aussichtsreichster Kandidat), alternative Exponenten
(unwahrscheinlich: $p_3^{\mathrm{PDG}} = 1{,}8081$ liegt nur $0{,}45\,\%$ über $9/5$).
$\Delta m_{32}^2$ bleibt \textbf{[S]}.
Prüfskript: \texttt{320\_dm32\_neutrino\_verify.py} (10 Assertionen).\\
\emph{Deklariert 23.~August 2026}

\medskip

\textbf{R88 -- $n_{\mathrm{thresh}}$ ist skalenabhängig; Matzkes Skala rekonstruiert (August 2026):}
Dok.~329 entscheidet die in Dok.~326 als [S] geführte Frage numerisch.
Die Kollapsbedingung Compton~$=$~Schwarzschild liefert $M_{\mathrm{coll}} = m_P/\sqrt{2}$ rein geometrisch und bitwertfrei \textbf{[B]}.
Mit der systemabhängigen Bit-Energie $E_{\mathrm{bit}}(L) = \hbar c/L$ (Dok.~257) folgt exakt $n_{\mathrm{thresh}}(L) = L/(\sqrt{2}\,\ell_P)$ -- linear in der Skala \textbf{[B]}.
Matzkes Wert $6{,}4097$ entspricht exakt $L = 2\pi\ell_P/\ln 2 = 9{,}0647\,\ell_P$ \textbf{[B]}; der „universelle" Bitwert ist die FFGFT-Bit-Energie bei dieser Skala.
Über die Skalen des Korpus variiert $n_{\mathrm{thresh}}$ um Faktor $4{,}6\times10^{29}$; eine Temperaturabweichung von $10\,\%$ kippt die Stabilitätsaussage für $\mathrm{Cl}(6)$.
Damit ist $n_{\mathrm{thresh}} = 6{,}41$ als universelle Zahl abgeschlossen negativ \textbf{[X]}.
Unberührt: Compton~$=$~Schwarzschild, $r_s$ ohne Feldgleichungen, $\mathrm{Cl}(6)$-Struktur \textbf{[B]}.
Prüfskript: \texttt{326\_nthresh\_skalenabhaengigkeit.py} (12 Assertionen).\\
\emph{Deklariert 24.~August 2026}

\medskip

\textbf{R89 -- Zeitfeld-Mediatormasse: eine Struktur, sektorabh\"angige Basismasse (August 2026):}
Der Korpus verwendet $m_T$ mit stark verschiedenen Zahlenwerten (Dok.~019/180/201: $m_T = \lambda/\xi$; Dok.~032: effektive Torsionsmasse $\approx 5{,}220$~GeV). Das ist kein Widerspruch: Beide folgen der Regel $m_T = M_\mathrm{Basis}/\xi$, woraus f\"ur die Reichweite des massiven Zeitfelds $\hbar/(m_T c) = \xi\cdot\lambda_C(M_\mathrm{Basis})$ folgt \textbf{[B]} --- universelle Form, sektorabh\"angiger Wert, strukturgleich mit $E_\mathrm{bit} = \hbar c/L$ (Dok.~257, 329).
Leptonsektor ($M = m_e$): $m_e/\xi = 3{,}83$~GeV, mit den $O(1)$-Faktoren aus Dok.~032 ($\sin\pi\xi$, $\pi^2$, $\sqrt{\alpha/K_\mathrm{frak}}$, $R_f$) $5{,}220$~GeV --- Reichweite $3{,}78\times10^{-17}$~m.
Fundamentaler Sektor ($M = m_P$): $m_T = m_P/\xi = 7500\,m_P$ --- Reichweite $2{,}155\times10^{-39}$~m $= L_0$ exakt.
Dok.~032 bezeichnet seine Gr\"o\ss{}e selbst als \emph{effektive} Torsionsmasse; die Namensgleichheit ist keine Gr\"o\ss{}engleichheit.
Pr\"ufskript: \texttt{180\_mT\_struktur\_L0\_verify.py} (17 Assertionen).\\
\emph{Deklariert 24.~August 2026}

\medskip

\noindent\textit{Einordnung nach Dok.~287:} Die Lagrangian-Beschreibung
($m_T = M_\mathrm{Basis}/\xi$) und die Bit-Energie-Beschreibung
($E_\mathrm{bit} = \hbar c/L$) sind \emph{verschiedene Darstellungen derselben
Boden-Struktur} -- nach Dok.~287 vermutlich eine \textbf{Bijektion}: verlustfrei
und gegenseitig \"ubersetzbar, aber keine vollst\"andiger als die andere.
Der generative Kern, der nichts verliert, ist der $T^4$-Torus (Dok.~287, Dok.~243);
alle anderen Beschreibungen stehen zu ihm in einer von drei Beziehungen:
Identit\"at, Bijektion oder Projektion.

\medskip
\textbf{R90 -- Der Exponent in $L_0 = \xi\cdot\ell_P$ ist hergeleitet (August 2026):}
Dok.~180 begr\"undet $L_0 = \xi\cdot\ell_P$ mit der $\xi$-Skalierung des Kopplungsterms $g_T^\ell = \xi m_\ell$ --- ein Analogieschluss.
Der Exponent folgt jedoch \emph{zwingend} aus dem Massenterm des Lagrangians: $m_T = M_\mathrm{Basis}/\xi^1$ und $\lambda_C \sim 1/m$ ergeben Reichweite $\sim\xi^1$ \textbf{[B]} (R89).
Gegenprobe: $m_T \sim \xi^{-n}$ erg\"abe $L_0 = \xi^n\ell_P$; der Exponent der Reichweite ist der negative Exponent von $m_T$.
$L_0$ ist damit die Compton-Reichweite des Zeitfeld-Mediators im fundamentalen Sektor --- unterhalb davon kann das Zeitfeld nicht mehr vermitteln, weshalb \enquote{Abstand} seine Bedeutung verliert. Das erkl\"art zugleich, \emph{warum} \"uberhaupt ein Boden existiert: weil der Mediator massiv ist.
Offen bleibt allein die Wahl $M_\mathrm{Basis} = m_P$ f\"ur den fundamentalen Sektor; sie f\"uhrt keinen neuen Parameter ein ($\lambda = 1$ in Planck-Einheiten), ist aber nicht unabh\"angig begr\"undet [S].
Ohne Folgen f\"ur Dok.~327 (R86): Beschr\"anktheit verlangt nur einen endlichen Boden \textbf{[B]}.\\
\emph{Deklariert 24.~August 2026}

\medskip
\textbf{R91 -- $\lambda$ in $m_T = \lambda/\xi$ ist eine Basismasse (August 2026):}
Dok.~019 und 201 bezeichnen $\lambda$ als \enquote{Higgs-Kopplungsparameter}. Das ist irref\"uhrend:
(i) Dimensional muss $\lambda$ eine \emph{Masse} sein, da $m_T$ eine Masse und $\xi$ dimensionslos ist \textbf{[B]}; der Higgs-Quartic $\lambda_h \approx 0{,}129$ ist dimensionslos.
(ii) $\lambda_h$ l\"auft unter der RGE und ist ein Niederenergie-Parameter der elektroschwachen Skala; ein dort gemessener Wert ist an der Planck-Skala nicht einsetzbar.
(iii) Als dimensionsloser Faktor in Planck-Einheiten interpretiert erg\"abe $\lambda_h$ die Reichweite $7{,}75\cdot L_0$ --- verfehlt $L_0$.
Korrekt ist $\lambda = M_\mathrm{Basis}$ des jeweiligen Sektors, im fundamentalen Fall $m_P$ (R89/R90).
Eine Higgs-Verbindung kann allenfalls eine effektive Aussage im elektroschwachen Sektor sein, nie die Definition. Die Beschriftung in Dok.~019/201 ist zu korrigieren [S].\\
\emph{Deklariert 24.~August 2026}

\medskip
\textbf{R92 -- Zwei Zuordnungsfehler zum Zeitfeld (August 2026):}
(i) \emph{Dok.~250}: Der Satz \enquote{Der Term $-\tfrac12 m_T^2(\partial m)^2$ mit $m_T = \xi/\ell$ (Dok.~201)} enth\"alt zwei Fehler.
In nat\"urlichen Einheiten ist $m = 1/\ell$, also $\xi/\ell = \xi\cdot m$ --- das ist die \emph{Kopplung} $g_T^\ell = \xi m_\ell$, nicht die Mediatormasse $m_T = M_\mathrm{Basis}/\xi$ \textbf{[B]}.
Ferner lautet der Massenterm $-\tfrac12 m_T^2(\Delta m)^2$, nicht $(\partial m)^2$; letzteres w\"are ein Beitrag zum kinetischen Term.
Der physikalische Punkt (der Massenterm liefert die Steifigkeit gegen Divergenz von $\Delta m$) bleibt richtig; nur die Formel ist falsch zitiert [S].
(ii) \emph{Dok.~190, Pr\"azisierung~4}: Die als interne Konsistenzpr\"ufung gef\"uhrte Aussage, ein Objekt mit $M_0 = \xi m_P/2$ h\"atte den Schwarzschild-Radius genau $L_0$, ist eine \emph{Identit\"at}: aus $r_s = 2GM/c^2$ und $\ell_P = G m_P/c^2$ folgt $r_s(x\,m_P/2) = x\,\ell_P$ f\"ur jedes $x$ und jeden Exponenten \textbf{[B]}.
Sie testet weder $\xi$ noch den Exponenten in $L_0$; die Einstufung als \enquote{Probe} ist zu hoch gegriffen. Nichttrivial bliebe allein die dort angef\"ugte Bedingung, ob $M_0$ mit einer Skalenstufe der $\xi$-Hierarchie zusammenf\"allt --- offen [S].\\
\emph{Deklariert 24.~August 2026}

\medskip
